% SOURCE_AUTHORITY: sga5_fr_workpass.tex, lines 7533--7576 (LF-normalized unit).
% SOURCE_SHA256: 791F4EFFC5E02832D5D77ED03518C8156D6F07E4C8238B03545DB93D883FBB28
\begin{prooftext}[Demonstração de 4.2.7]
Vamos retomar a notação de (4.2.3). A questão é ver que o morfismo
\[
\omega_{pq}:J^p/J^{p+1}\otimes_A F^q(X_q)\longrightarrow F^{p+q}(X_{p+q})
\]
envia
\[
J^p/J^{p+1}\otimes_A
\bigl(J^{q-s}X_q\cap F^q(X_q)\bigr)
\]
em
\[
J^{p+q-s}X_{p+q}\cap F^{p+q}(X_{p+q}).
\]
Ora, o morfismo canônico
\[
\rho:X_{p+q}/F^{q+1}(X_{p+q})\longrightarrow X_q
\]
é um isomorfismo e envia
\[
\bigl(J^{q-s}X_{p+q}+F^{q+1}(X_{p+q})\bigr)/F^{q+1}(X_{p+q})
\]
sobre \(J^{q-s}X_q\). Disso se deduz que a imagem inversa por
\(\id(J^p/J^{p+1})\otimes_A\varphi\) (4.2.3) de
\[
J^p/J^{p+1}\otimes_A
\bigl(J^{q-s}X_q\cap F^q(X_q)\bigr)
\]
é
\[
J^p/J^{p+1}\otimes_A
[(J^{q-s}X_{p+q}\cap F^q(X_{p+q}))+F^{q+1}(X_{p+q})]/F^{q+1}(X_{p+q}).
\]
Agora está claro que a imagem por \(\beta\) (4.2.3) deste último objeto está contida em
\[
J^{p+q-s}X_{p+q}\cap J^pF^q(X_{p+q}),
\]
e, portanto, com maior razão, em
\[
J^{p+q-s}X_{p+q}\cap F^{p+q}(X_{p+q}).
\]
Daí segue o lema.
\end{prooftext}
